\section{Implementation Details}\label{sec:implementation_details}
\subsection[ ]{Handling Conditional Implementations}

\begin{frame}{\secname}
	\framesubtitle{\subsecname}

	Conditions defined in the \inlinerust{when} macro are normalized into \alert{DNF}, % Disjunctive Normal Form (disjunction of conjunctions)
	flattening when possible, e.g.:
	% Once we have the dnf conditions, each disjunct condition is treated as a separate implementation
	% This allows to handle each disjunction independently during method resolution

	\smallskip

	{\footnotesize\inlinerust{all(T = Vec<U>, any(U = i32, U = f64))} $\rightarrow$ \inlinerust{any(T = Vec<i32>, T = Vec<f64>)}}

	\bigskip
	\pause

	A \alert{partial order} defined over the predicates is used to determine the specificity of each implementation, e.g.:

	\smallskip

	{\footnotesize\inlinerust{T = &i32} $>$ \inlinerust{T = &U} $>$ \inlinerust{T: Copy}}

	\bigskip
	\pause

	% In addition to the two macros
	We need a way to collect all the traits defined in the crate % to be used in the monomorphization phase
	and to handle the order of their implementations. % such that a method call can also call an implementation defined later in the code
	In Rust, this is done with \alert{Build Scripts}, e.g.:

	\smallskip

	{\footnotesize\inlinerust{mono_spec::handle_order();}}

\end{frame}